%-----------------------------------------------------------------------------------------------------------------------------------------------%
%	The MIT License (MIT)
%
%	Copyright (c) 2015 Jan Küster
%
%	Permission is hereby granted, free of charge, to any person obtaining a copy
%	of this software and associated documentation files (the "Software"), to deal
%	in the Software without restriction, including without limitation the rights
%	to use, copy, modify, merge, publish, distribute, sublicense, and/or sell
%	copies of the Software, and to permit persons to whom the Software is
%	furnished to do so, subject to the following conditions:
%	
%	THE SOFTWARE IS PROVIDED "AS IS", WITHOUT WARRANTY OF ANY KIND, EXPRESS OR
%	IMPLIED, INCLUDING BUT NOT LIMITED TO THE WARRANTIES OF MERCHANTABILITY,
%	FITNESS FOR A PARTICULAR PURPOSE AND NONINFRINGEMENT. IN NO EVENT SHALL THE
%	AUTHORS OR COPYRIGHT HOLDERS BE LIABLE FOR ANY CLAIM, DAMAGES OR OTHER
%	LIABILITY, WHETHER IN AN ACTION OF CONTRACT, TORT OR OTHERWISE, ARISING FROM,
%	OUT OF OR IN CONNECTION WITH THE SOFTWARE OR THE USE OR OTHER DEALINGS IN
%	THE SOFTWARE.
%	
%
%-----------------------------------------------------------------------------------------------------------------------------------------------%


%============================================================================%
%
%	DOCUMENT DEFINITION
%
%============================================================================%

%we use article class because we want to fully customize the page and dont use a cv template
\documentclass[10pt,A4]{article}	


%----------------------------------------------------------------------------------------
%	ENCODING
%----------------------------------------------------------------------------------------

%we use utf8 since we want to build from any machine
\usepackage[utf8]{inputenc}		

\usepackage{hyperref}

\makeatletter
\hypersetup{
	%pagebackref=true,
	colorlinks=true,       		% false: boxed links; true: colored links
linkcolor=black,          	% color of internal links
	citecolor=black,        		% color of links to bibliography
	filecolor=magenta,      		% color of file links
urlcolor=blue,
	bookmarksdepth=4
}
\makeatother
%----------------------------------------------------------------------------------------
%	LOGIC
%----------------------------------------------------------------------------------------

% provides \isempty test
\usepackage{xifthen}

%----------------------------------------------------------------------------------------
%	FONT
%----------------------------------------------------------------------------------------

% some tex-live fonts - choose your own

%\usepackage[defaultsans]{droidsans}
%\usepackage[default]{comfortaa}
%\usepackage{cmbright}
\usepackage[default]{raleway}
%\usepackage{fetamont}
%\usepackage[default]{gillius}
%\usepackage[light,math]{iwona}
%\usepackage[thin]{roboto} 

% set font default
\renewcommand*\familydefault{\sfdefault} 	
\usepackage[T1]{fontenc}

% more font size definitions
\usepackage{moresize}		


%----------------------------------------------------------------------------------------
%	PAGE LAYOUT  DEFINITIONS
%----------------------------------------------------------------------------------------

%debug page outer frames
%\usepackage{showframe}			


%define page styles using geometry
\usepackage[a4paper]{geometry}		

% for example, change the margins to 2 inches all round
\geometry{top=1cm, bottom=1cm, left=1.5cm, right=1.5cm} 	

%use customized header
\usepackage{fancyhdr}				
\pagestyle{fancy}

%less space between header and content
\setlength{\headheight}{-5pt}		


%customize entries left, center and right
%\lhead{}
%\chead{}
%\rhead{}


%indentation is zero
\setlength{\parindent}{0mm}

%----------------------------------------------------------------------------------------
%	TABLE /ARRAY DEFINITIONS
%---------------------------------------------------------------------------------------- 

%for layouting tables
\usepackage{multicol}			
\usepackage{multirow}

%extended aligning of tabular cells
\usepackage{array}

\newcolumntype{x}[1]{%
>{\raggedleft\hspace{0pt}}p{#1}}%


%----------------------------------------------------------------------------------------
%	GRAPHICS DEFINITIONS
%---------------------------------------------------------------------------------------- 

%for header image
\usepackage{graphicx}

%for floating figures
\usepackage{wrapfig}
\usepackage{float}
%\floatstyle{boxed} 
%\restylefloat{figure}

%for drawing graphics		
\usepackage{tikz}				
\usetikzlibrary{shapes, backgrounds,mindmap, trees}


%----------------------------------------------------------------------------------------
%	Color DEFINITIONS
%---------------------------------------------------------------------------------------- 

\usepackage{color}

%accent color
\definecolor{sectcol}{RGB}{0,150,255}

%dark background color
\definecolor{bgcol}{RGB}{110,110,110}

%light background / accent color
\definecolor{softcol}{RGB}{225,225,225}


%============================================================================%
%
%
%	DEFINITIONS
%
%
%============================================================================%

%----------------------------------------------------------------------------------------
% 	HEADER
%----------------------------------------------------------------------------------------

% remove top header line
\renewcommand{\headrulewidth}{0pt} 

%remove botttom header line
\renewcommand{\footrulewidth}{0pt}	  	

%remove pagenum
\renewcommand{\thepage}{}	

%remove section num		
\renewcommand{\thesection}{}			

%----------------------------------------------------------------------------------------
% 	ARROW GRAPHICS in Tikz
%----------------------------------------------------------------------------------------

% a six pointed arrow poiting to the left
\newcommand{\tzlarrow}{(0,0) -- (0.2,0) -- (0.3,0.2) -- (0.2,0.4) -- (0,0.4) -- (0.1,0.2) -- cycle;}	

% include the left arrow into a tikz picture
% param1: fill color
%
\newcommand{\larrow}[1]
{\begin{tikzpicture}[scale=0.58]
	 \filldraw[fill=#1!100,draw=#1!100!black]  \tzlarrow
 \end{tikzpicture}
}

% a six pointed arrow poiting to the right
\newcommand{\tzrarrow}{ (0,0.2) -- (0.1,0) -- (0.3,0) -- (0.2,0.2) -- (0.3,0.4) -- (0.1,0.4) -- cycle;}

% include the right arrow into a tikz picture
% param1: fill color
%
\newcommand{\rarrow}
{\begin{tikzpicture}[scale=0.7]
	\filldraw[fill=sectcol!100,draw=sectcol!100!black] \tzrarrow
 \end{tikzpicture}
}



%----------------------------------------------------------------------------------------
%	custom sections
%----------------------------------------------------------------------------------------

% create a coloured box with arrow and title as cv section headline
% param 1: section title
%

\newcommand{\eventspace}{
	\vspace{0.05cm}
}

\newcommand{\sectionspace}{
	\vspace{0.1cm}
}

\newcommand{\cvsection}[1]
{
	\begin{center}
		\large\textcolor{sectcol}{\textbf{#1}}
	\end{center}
	\sectionspace
}

%create a coloured arrow with title as cv meta section section
% param 1: meta section title
%
\newcommand{\metasection}[2]
{
%\begin{tabular*}{1\textwidth}{r r}
\footnotesize{#2} \hspace*{\fill} \footnotesize{#1}\\[1pt]
%\end{tabular*}
}

%----------------------------------------------------------------------------------------
%	 CV EVENT
%----------------------------------------------------------------------------------------

% creates a stretched box as cv entry headline followed by some paragraphs about 
% the work you did
% param 1:	event time i.e. 2014 or 2011-2014 etc.
% param 2:	event name (what did you do?)
% param 3:	institution (where did you work / study)
% param 4:	list of paragraphs
%
\newcommand{\cvevent}[4]
{

\begin{tabular*}{1\textwidth}{p{13cm}  x{4.5cm}}
	\textbf{#2} - \textcolor{bgcol}{#3} &   \vspace{2.5pt}\textcolor{sectcol}{#1}
\end{tabular*}

\vspace{-8pt}
\textcolor{softcol}{\hrule}
\vspace{6pt}

	\foreach \desc in {#4}{
		$\cdot$ \desc\\[3pt]
	}
	
\eventspace
}

\newcommand{\cveventextra}[4]
{

\begin{tabular*}{1\textwidth}{p{13cm}  x{4.5cm}}
	\textbf{#2} - \textcolor{bgcol}{#3} &   \vspace{2.5pt}\textcolor{sectcol}{#1}
\end{tabular*}

\vspace{-4pt}
\textcolor{softcol}{\hrule}
\vspace{6pt}

	\foreach \desc in {#4}{
		$\cdot$ \desc\\[3pt]
	}
	
\eventspace
}

% creates a stretched box as 
\newcommand{\cveventmeta}[2]
{
	\mbox{\mystrut \hspace{87pt}\textit{#1}}\\
	#2
}

%----------------------------------------------------------------------------------------
% CUSTOM STRUT FOR EMPTY BOXES
%----------------------------------------- -----------------------------------------------
\newcommand{\mystrut}{\rule[-.3\baselineskip]{0pt}{\baselineskip}}

%----------------------------------------------------------------------------------------
% CUSTOM LOREM IPSUM
%----------------------------------------------------------------------------------------
\newcommand{\lorem}
{Lorem ipsum dolor sit amet, consectetur adipiscing elit. Donec a diam lectus.}



%============================================================================%
%
%
%
%	DOCUMENT CONTENT
%
%
%
%============================================================================%
\begin{document}


%use our custom fancy header definitions
\pagestyle{fancy}	





%----------------------------------------------------------------------------------------
%	HEADER IMAGE
%----------------------------------------------------------------------------------------

%\begin{figure}[H]
%\begin{flushright}
%	\includegraphics[trim= 320 130 460 210,clip,width=0.2\linewidth]{myfoto.jpg}	%trimming relative to image size!
%\end{flushright}
%\end{figure}


%---------------------------------------------------------------------------------------
%	TITLE HEADLINE
%----------------------------------------------------------------------------------------
\vspace{-8pt}
\begin{center}
	\HUGE \textsc{Paulo Fylippe Sell} \\[2pt]
\end{center}



\vspace{6pt}


%---------------------------------------------------------------------------------------
%	META SECTION
%----------------------------------------------------------------------------------------
\metasection{São José, Santa Catarina}{\textbf{Status:} Engenheiro de telecomunicações}
\metasection{\href{https://paulosell.github.io}{paulosell.github.io}}{\textbf{Áreas:} 
Sistemas embarcados, Linux, Acelerador em hardware}
\metasection{sellpaulof@gmail.com}{\textbf{Techs:} C/C++, Python, VHDL, Git}
\metasection{+55 48 984053102}{\textbf{Links:} \href{https://github.com/paulosell}{github.com/paulosell} $\cdot$ \href{https://www.linkedin.com/in/paulosell/}{linkedin.com/in/paulosell/}} 
\vspace{-2pt}
\textcolor{softcol}{\hrule}
\vspace{6pt}

\normalsize

%---------------------------------------------------------------------------------------
%	SUMMARAY (optional)
%----------------------------------------------------------------------------------------
\vspace{-6pt}
\cvsection{Resumo}
Engenheiro de Telecomunicações morador da cidade de São José, Santa Catarina. Formado no Instituto Federal de Santa Catarina, com tese de graduação relacionada com o desenvolvimento de um framework para atualizações seguras e confiáveis de firmware para dispositivos embarcados baseados em microcontroladores. Atualmente cursando mestrado Computação Aplicada na Universidade do Vale do Itajaí (UNIVALI), com especialização em Sistemas Embarcados e pesquisa com objetivo de melhorar a eficiência energética e aumentar a taxa de transferência em aceleradores de hardware baseados em FPGA para redes neurais profundas.\\

%============================================================================%
%
%	CV SECTIONS AND EVENTS (MAIN CONTENT)
%
%============================================================================%

%---------------------------------------------------------------------------------------
%	EXPERIENCE
%----------------------------------------------------------------------------------------
\cvsection{Experiências profissionais}

\cvevent{ago 2023 - presente}{Engenheiro de software embarcado}{GE Vernova}{
	{Atuação no desenvolvimento de software embarcado para dispositivos de controle, monitoramento e automação de estações de energia elétrica, incluindo dispositivos como Digital Fault Recorders (DFR) e Merging Units (MU).},
	{Virtualização de dispositivos legado utilizando Xen Hypervisor.},
	{Desenvolvimento com C, C++ e Shell script para criação de soluções em sistemas Linux embarcados e integração com hardware e FPGA.},
	{Contribuição para aprimoramento da estabilidade e desempenho dos sistemas embarcados por meio do desenvolvimento de drivers e kernel do Linux.}
}

\cvevent{fev 2022 - jul 2023}{Analista de desenvolvimento de produto}{Intelbras}{
{Responsável pelo desenvolvimento de software embarcado, incluindo a implementação de novos recursos de rede, gerenciamento de comportamentos do dispositivo e fornecendo informações de status em tempo real por meio de APIs REST com HTTP e broker MQTT para aproximadamente 1 milhão de roteadores home \& office baseados em Linux com um time de 4 pessoas.},
{Desenvolvimento com C e Shell script para soluções de hardware entregues por fornecedores baseadas SDK com Linux embarcado, modificando drivers e kernel do Linux e aplicações relacionadas ao gerenciamento dos produtos.},
{Desenvolvimento de produtos baseados em microcontroladores ARM para soluções Alexa Voice Service (AVS) e ecossistema Alexa/AWS (IAM, S3, IoT Core) com FreeRTOS.}
}

\cvevent{maio 2021 - fev 2022}{Desenvolvedor de software / Pesquisador}{Instituto Federal de Santa Catarina}{
{Atuação no grupo de pesquisa LaSED, contribuindo com o desenvolvimento da especificação técnica de requisitos de um novo dispositivo embarcado microcontrolado de autorização fiscal \href{https://www.sef.sc.gov.br/arquivos_portal/servicos/159/Especificacao_de_Requisitos_do_DAF___versao_3.0.0.pdf}{(Dispositivo Autorizador Fiscal - DAF)}, a ser utilizado na autenticação de Nota Fiscal Consumidor Eletrônica (NFC-e) em todo o Estado de Santa Catarina, de acordo com a legislação brasileira.},
{Descrição de requisitos de arquitetura de software e hardware, de protocolo de comunicação e de hardware seguro; descrição de casos de uso do dispositivo; escrita da API e protocolo de comunicação do DAF; desenvolvimento do firmware da prova de conceito do DAF; integração do DAF com softwares complementares. Atuação também na elaboração do conjunto de instruções e regras para homologações de DAF para o mercado.},
{Desenvolvimento de provas de conceito com C++, Java, Python.},
{Experiências com stacks como: Flask, Docker, Redis, MySQL, Protocolo USB, criptografia de chave pública com JWT, hardware seguro.}
}
%
%---------------------------------------------------------------------------------------
%	EDUCATION SECTION
%--------------------------------------------------------------------------------------
\newpage
\cvsection{Educação}

\cvevent{2023 - presente}{Mestrado em computação aplicada}{Universidade do Vale do Itajaí }{
	{Especialização em sistemas embarcados.},
	{Melhoria da eficiência energética e aumento da taxa de transferência em aceleradores de hardware baseados em FPGA para redes neurais profundas.}
}

\cveventextra{2015 - 2021}{Graduação em engenharia de telecomunicações}{Instituto Federal de Santa Catarina}{
{Desenvolvimento de um framework para atualização segura e confiável de firmware para dispositivos embarcados baseados em microcontroladores utilizando ARM TrustZone.},
	{
		Prêmio de \href{https://sol.sbc.org.br/index.php/sbseg_estendido/article/view/17354/17192}{melhor artigo completo} no 15º Workshop de Iniciação Científica e Pesquisa de Graduação (WTICG) do 21º Simpósio Brasileiro em Segurança da Informação e Sistemas Computacionais (SBSeg 2021).
	}
}

% \cvevent{2011 - 2014}{Vocational High School}{Federal Institute of Santa Catarina}{
% 	{Telecommunications technician}
% }

\cvsection{Experiências acadêmicas}

\cvevent{2020 - 2021}{Projeto de pesquisa}{Instituto Federal de Santa Catarina}{
	{Desenvolvimento da especificação e do protótipo de um dispositivo embarcado de autorização fiscal seguro de baixo custo para o estado de Santa Catarina.}
}

\cvevent{2019 - 2019}{Projeto de pesquisa}{Instituto Federal de Santa Catarina}{
	{Desenvolvimento de um bloco de processamento de sinal digital (DSP) para uma Transformada Cosseno Discreta Modificada (MDCT).},
	{Uso de Matlab e VHDL para o desenvolvimento do bloco e sintetização em hardware.}}



\cvsection{Áreas de conhecimento}

$\cdot$  Sistemas embarcados, comunicação sem fio, projeto de protocolos, redes neurais profundas, RF, circuitos elétricos, sistemas de telecomunicações, sistemas distribuídos.


%\textcolor{softcol}{\hrule}

%

%--------------------------------------------------------------------------------------------------
%	ARTIFICIAL FOOTER (fancy footer cannot exceed linewidth) 
%--------------------------------------------------------------------------------------------------

%============================================================================%
%
%
%
%	DOCUMENT END
%
%
%
%============================================================================%
\end{document}
